% संपूर्ण मराठी ग्रंथातील प्रकरण 19: प्रतिमान उपपत्तीची मूलतत्त्वे
% हा संपादनयोग्य स्रोतखंड आहे. संकलनासाठी संपूर्ण स्रोतसंग्रहातील
% अक्षररूपे, पूर्वभाग, चिन्हव्याख्या आणि आकृती-स्रोत आवश्यक आहेत.
% मूळ: Open Logic Project; मजकूर आणि रूपांतर CC BY 4.0.
% यंत्रानुवाद, दुरुस्ती आणि तपासणी: OpenAI Codex —
% GPT-5.6 Sol आणि GPT-6 Sol, दोन्ही Ultra effort.
% दुरुस्ती-पुनरावलोकन: GPT-6.1 Sol, Ultra effort.
\chapter{प्रतिमान उपपत्तीची मूलतत्त्वे}\label{mod:bas::chap}








\section{न्यूनीकृत आणि विस्तारित रचना}\label{mod:bas:red:sec}

ज्या भाषांत काही चिन्हे समान आहेत अशा भाषांची, तसेच त्या भाषांसाठीच्या
संरचनेची, तुलना करणे अनेकदा उपयुक्त किंवा आवश्यक असते. सर्वांत
नेहमीचे प्रकरण असे असते की भाषा~$\Lang{L}$ मधील सर्व चिन्हे
भाषा~$\Lang{L'}$ मध्येही असतात, म्हणजे $\Lang{L} \subseteq
\Lang{L'}$. मग अतिरिक्त चिन्हांची अर्थनिर्धारणे जोडून आणि समान चिन्हांची
अर्थनिर्धारणे तशीच ठेवून $\Lang{L}$-संरचना~$\Struct{M}$ चा विस्तार
नेहमीच $\Lang{L'}$-संरचनांमध्ये करता येतो. दुसरीकडे,
$\Lang{L}$ मध्ये नसलेल्या चिन्हांची अर्थनिर्धारणे केवळ ``विसरून''
$\Lang{L'}$-रचना~$\Struct{M'}$ पासून $\Lang{L}$-रचना मिळवता येते.

\begin{defn}
\label{mod:bas:red:defn:reduct}
समजा $\Lang L \subseteq \Lang L'$, $\Struct M$ ही
$\Lang L$-संरचना आणि $\Struct M'$ ही $\Lang L'$-संरचना आहे.
$\Struct M$ ही $\Struct M'$ ची $\Lang L$ पर्यंतची \emph{न्यूनीकृत रचना}
आणि $\Struct M'$ ही $\Struct M$ ची $\Lang L'$ पर्यंतची
\emph{विस्तारित रचना} असेल तेव्हा आणि केवळ तेव्हाच पुढील अटी पूर्ण होतात:
\begin{enumerate}
\item $\Domain{M} = \Domain{M'}$
\item प्रत्येक स्थिरांक~$c \in \Lang L$ साठी $\Assign{c}{M} =
  \Assign{c}{M'}$.
\item प्रत्येक फलन~$f \in \Lang L$ साठी $\Assign{f}{M} =
  \Assign{f}{M'}$.
\item प्रत्येक विधेय~$P \in \Lang L$ साठी $\Assign{P}{M} =
  \Assign{P}{M'}$.
\end{enumerate}
\end{defn}

\begin{prop}
\label{mod:bas:red:prop:reduct}
एखादी $\Lang{L}$-संरचना~$\Struct{M}$ ही एखाद्या
$\Lang{L'}$-संरचना~$\Struct{M'}$ ची न्यूनीकृत रचना असेल, तर सर्व
$\Lang{L}$-वाक्ये~$\varphi$ साठी,
\[\Sat{M}{\varphi} \text{ तेव्हा आणि केवळ तेव्हाच } \Sat{M'}{\varphi}.
\]
\end{prop}

\begin{proof}
  सराव.
\end{proof}

\begin{prob}
\ref{mod:bas:red:prop:reduct} सिद्ध करा.
\end{prob}

\begin{defn}
आपल्याकडे $\Lang{L}$-रचना $\Struct{M}$ असेल, एका $n$-स्थानी
विधेय~$P$ ची भर घालून मिळालेला $\Lang{L}$ चा विस्तार
$\Lang{L'} = \Lang{L} \cup \{P\}$ असेल, आणि $R \subseteq \Domain{M}^n$
हा $n$-स्थानी संबंध असेल, तर $\Assign{P}{M'} = R$ असलेली
$\Struct{M}$ ची विस्तारित रचना~$\Struct{M'}$ दर्शवण्यासाठी आपण
$\Expan{M}{R}$ असे लिहितो.
\end{defn}











\section{उपसंरचना}\label{mod:bas:sub:sec}

एखाद्या संरचना~$\Struct{M}$ चा प्रांत हा दुसऱ्या
$\Struct{M'}$ चा उपसंच असू शकतो. पण $\Struct{M}$ हा $\Struct{M'}$ चा
``भाग'' आहे असे आपण अर्थातच फक्त तेव्हाच मानावे, जेव्हा $\Domain{M}
\subseteq \Domain{M'}$ असण्याबरोबरच भाषेतील चिन्हांचे अर्थनिर्धारण करताना
$\Struct{M}$ आणि $\Struct{M'}$ किमान सामाईक भाग~$\Domain{M}$ वर
``सहमत'' असतात.

\begin{defn}
\label{mod:bas:sub:defn:substructure}
त्याच भाषेसाठी~$\Lang L$ संरचना $\Struct M$ आणि
$\Struct M'$ दिल्या असता, $\Struct M$ ही $\Struct M'$ ची
\emph{उपसंरचना} आणि $\Struct M'$ ही $\Struct M$ ची
\emph{वर्धित रचना} आहे, असे आपण म्हणतो; $\Struct M \substruct \Struct M'$
असे लिहितो; आणि तेव्हा आणि केवळ तेव्हाच पुढील अटी पूर्ण होतात:
\begin{enumerate}
\item $\Domain{M} \subseteq \Domain{M'}$,
\item प्रत्येक स्थिर $c \in \Lang L$ साठी $\Assign{c}{M} =
    \Assign{c}{M'}$;
\item प्रत्येक $n$-स्थानी फलन $f \in \Lang L$ साठी
  सर्व $a_1$, \dots, $a_n \in \Domain{M}$ यांकरिता
  $\Assign{f}{M}(a_1, \dots, a_n) = \Assign{f}{M'}(a_1, \dots, a_n)$.
\item प्रत्येक $n$-स्थानी विधेय $R \in \Lang L$ साठी, सर्व
  $a_1$, \dots, $a_n \in \Domain{M}$ यांकरिता $\langle
  a_1, \dots, a_n\rangle \in \Assign{R}{M}$ तेव्हा आणि केवळ तेव्हाच,
  जेव्हा $\langle a_1, \dots, a_n\rangle \in \Assign{R}{M'}$.
\end{enumerate}
\end{defn}

\begin{rem}
\label{mod:bas:sub:rem:substructure}
भाषेत एकही स्थिर किंवा फलने नसतील, तर कोणताही $N
\subseteq \Domain{M}$ हा $\Assign{R}{N} = \Assign{R}{M} \cap N^n$
असे ठेवून $\Struct M$ ची, प्रांत~$\Domain{N} = N$ असलेली,
उपसंरचना~$\Struct{N}$ निश्चित करतो.
\end{rem}













\section{सांत मर्यादेपलीकडील प्रसरण}\label{mod:bas:ove:sec}

\begin{thm}
\label{mod:bas:ove:overspill} वाक्यांचा एखादा संच $\Gamma$ याला हवे तितके मोठे
सांत प्रतिमान असतील, तर त्याला एक अनंत प्रतिमान असते.
\end{thm}

\begin{proof}
$c_0$, $c_1$, \dots ही गणनीय इतकी नवीन स्थिरे जोडून $\Gamma$ च्या भाषेचा
विस्तार करा आणि $\Gamma \cup \{c_i \neq c_j : i \neq j\}$ हा संच विचारात
घ्या. $\Gamma$ ला हवे तितके मोठे सांत प्रतिमान आहेत, याचा अर्थ प्रत्येक
$m >0$ साठी असा $n\ge m$ असतो की $\Gamma$ ला आकारमान~$n$ चे प्रतिमान
असते. यावरून $\Gamma \cup \{c_i \neq c_j : i \neq j\}$ सांततः
पूर्ततायोग्य आहे. संहततेनुसार $\Gamma \cup \{c_i \neq c_j : i \neq j\}$
याला प्रतिमान $\Struct M$ असते; त्याचा प्रांत अनंत असलाच पाहिजे, कारण ते
$c_i \neq c_j$ या सर्व असमानतांची पूर्ति करते.
\end{proof}

\begin{prop}
\label{mod:bas:ove:inf-not-fo}
कोणत्याही प्रथम-क्रम भाषेत असे कोणतेही वाक्य $\varphi$ नाही की ते
संरचना~$\Struct M$ मध्ये तेव्हा आणि केवळ तेव्हाच सत्य असते, जेव्हा
त्या संरचनेचा प्रांत $\Domain{M}$ अनंत असतो.
\end{prop}

\begin{proof}
असे एखादे $\varphi$ असते, तर त्याचे नकरण $\lnot \varphi$ सर्व आणि फक्त सांत
संरचनांमध्ये सत्य असते. म्हणून त्याला हवे तितके मोठे सांत प्रतिमान
असतात, पण अनंत प्रतिमान नसते; हे \ref{mod:bas:ove:overspill} शी विसंगत आहे.
\end{proof}











\section{समरूपी रचना}\label{mod:bas:iso:sec}

प्रथम-क्रमाच्या संरचना दोन प्रकारे एकसारख्या असू शकतात. त्यांच्या
सारखेपणाचा एक प्रकार असा की त्या तीच वाक्ये सत्य ठरवतात. अशा
संरचना ना आपण \emph{प्राथमिक सममूल्य} म्हणतो. परंतु रचना अतिशय
भिन्न असूनही तीच वाक्ये सत्य ठरवू शकतात---उदाहरणार्थ, एक रचना
गणनीय असू शकते आणि दुसरी नसू शकते. याचे कारण असे की
संरचनेची पुष्कळ वैशिष्ट्ये प्रथम-क्रम भाषांमध्ये व्यक्त करता येत
नाहीत: कधी भाषा पुरेशी समृद्ध नसते, तर कधी लोव्हेनहाइम--स्कोलेम
प्रमेयासारख्या प्रथम-क्रम तर्कशास्त्राच्या मूलभूत मर्यादा आड येतात. म्हणून
संरचनेच्या सारखेपणाचा दुसरा, अधिक काटेकोर प्रकार असा की त्या
मूलतः एकच असतात: त्यांना घडवणाऱ्या वस्तू वेगळ्या असतात, पण त्यांची
रचनात्मक वैशिष्ट्ये वेगळी नसतात. \emph{समरूपते} च्या संकल्पनेने हे नेमके
मांडता येते.

\begin{defn}
\label{mod:bas:iso:defn:elem-equiv}
एकाच भाषा~$\Lang{L}$ साठीच्या दोन संरचना $\Struct{M}$ आणि
$\Struct M'$ दिलेल्या असताना, $\Lang{L}$ मधील प्रत्येक वाक्य~$\varphi$
साठी $\Sat{M}{\varphi}$ तेव्हा आणि केवळ तेव्हाच जेव्हा $\Sat{M'}{\varphi}$ असेल,
तर $\Struct{M}$ ही $\Struct M'$ शी \emph{प्राथमिक सममूल्य} आहे असे म्हणतो
आणि $\Struct{M} \equiv \Struct M'$ असे लिहितो.
\end{defn}

\begin{defn}
\label{mod:bas:iso:defn:isomorphism} एकाच भाषा~$\Lang L$ साठीच्या दोन
संरचना $\Struct{M}$ आणि $\Struct M'$ दिलेल्या असताना, पुढील अटी
पूर्ण करणारे $h \colon \Domain{M} \to \Domain{M'}$ फलन असेल, तर
$\Struct{M}$ ही~$\Struct M'$ शी \emph{समरूपी} आहे असे म्हणतो आणि
$\Struct{M} \simeq \Struct M'$ असे लिहितो:
\begin{enumerate}
\item $h$ हे एकास-एक आहे: $h(x) =
  h(y)$ असेल, तर $x = y$;
\item $h$ हे आच्छादक आहे: प्रत्येक $y \in \Domain{M'}$ साठी
  $h(x) = y$ असा $x \in \Domain{M}$ असतो;
\item \label{mod:bas:iso:defn:iso-const}प्रत्येक स्थिरांक $c$ साठी:
  $h(\Assign{c}{M}) = \Assign{c}{M'}$;
\item \label{mod:bas:iso:defn:iso-pred}प्रत्येक $n$-स्थानी विधेय~$P$ साठी:
  \[  \tuple{a_1, \dots, a_n}\in \Assign{P}{M} \quad\text{तेव्हा आणि केवळ तेव्हाच जेव्हा}\quad
  \tuple{h(a_1), \dots, h(a_n)} \in \Assign{P}{M'};
  \]
\item \label{mod:bas:iso:defn:iso-func}प्रत्येक $n$-स्थानी फलन $f$ साठी:
  \[  h(\Assign{f}{M}(a_1, \dots, a_n)) =
  \Assign{f}{M'}(h(a_1), \dots, h(a_n)).
  \]
\end{enumerate}
\end{defn}

\begin{thm}
\label{mod:bas:iso:thm:isom}
$\Struct{M} \iso \Struct M'$ असेल, तर $\Struct{M} \elemequiv
\Struct{M'}$.
\end{thm}

\begin{proof}
$h$ हे $\Struct{M}$ पासून $\Struct M'$ वरचे समरूपण असू द्या. कोणत्याही
मूल्यांकन~$s$ साठी $h \circ s$ हे $h$ आणि $s$ यांचे संयोजन असते, म्हणजे
ते $\Struct{M'}$ मधील असे मूल्यांकन आहे की $(h \circ s)(x) = h(s(x))$.
$t$ आणि $\varphi$ यांवरील विगमनाने पुढील अधिक सबळ दावे सिद्ध करता येतात:
\begin{enumerate}
  \item[a.] $h(\Value{t}{M}[s]) = \Value{t}{M'}[h\circ s]$.
  \item[b.] $\Sat{M}{\varphi}[s]$ तेव्हा आणि केवळ तेव्हाच जेव्हा $\Sat{M'}{\varphi}[h \circ s]$.
\end{enumerate}
पहिला दावा~$t$ च्या गुंतागुंतीवरील विगमनाने सिद्ध होतो.
\begin{enumerate}
\item $t \ident c$ असेल, तर $\Value{c}{M}[s] = \Assign{c}{M}$ आणि
  $\Value{c}{M'}[h \circ s] = \Assign{c}{M'}$. म्हणून
  $h(\Value{t}{M}[s]) = h(\Assign{c}{M}) = \Assign{c}{M'}$
  (\ref{mod:bas:iso:defn:isomorphism} मधील \ref{mod:bas:iso:defn:iso-const} नुसार) $=
  \Value{t}{M'}[h \circ s]$.
\item $t \ident x$ असेल, तर $\Value{x}{M}[s] = s(x)$ आणि
  $\Value{x}{M'}[h \circ s] = h(s(x))$. म्हणून $h(\Value{x}{M}[s]) =
  h(s(x)) = \Value{x}{M'}[h \circ s]$.
\item $t \ident f(t_1, \dots, t_n)$ असेल, तर
  \begin{align*}
    \Value{t}{M}[s] & = \Assign{f}{M}(\Value{t_1}{M}[s], \dots, \Value{t_n}{M}[s]) \quad\text{आणि}\\
  \Value{t}{M'}[h \circ s] & = \Assign{f}{M'}(\Value{t_1}{M'}[h \circ
    s], \dots, \Value{t_n}{M'}[h \circ s]).
  \end{align*}

  प्रत्येक $i$ साठी $h(\Value{t_i}{M}[s])
  = \Value{t_i}{M'}[h\circ s]$ हे विगमन गृहीतक आहे. म्हणून
  \begin{align}
    h(\Value{t}{M}[s])
    & = h(\Assign{f}{M}(\Value{t_1}{M}[s], \dots, \Value{t_n}{M}[s]) \notag\\
    & = \Assign{f}{M'}(h(\Value{t_1}{M}[s]), \dots,
    h(\Value{t_n}{M}[s])) \label{mod:bas:iso:iso-1}\\
    & = \Assign{f}{M'}(\Value{t_1}{M'}[h \circ s], \dots,
    \Value{t_n}{M'}[h \circ s]) \label{mod:bas:iso:iso-2}\\
    & = \Value{t}{M'}[h\circ s] \notag
  \end{align}
  येथे \ref{mod:bas:iso:iso-1} हे \ref{mod:bas:iso:defn:isomorphism} मधील
  \ref{mod:bas:iso:defn:iso-func} वरून, तर \ref{mod:bas:iso:iso-2} हे विगमन गृहीतकावरून मिळते.
\end{enumerate}
(b) भागाचा पुरावा सरावासाठी सोडला आहे.

$\varphi$ हे वाक्य असेल, तर मूल्यांकने~$s$ आणि $h \circ s$ महत्त्वाची राहत
नाहीत; म्हणून $\Sat{M}{\varphi}$ तेव्हा आणि केवळ तेव्हाच जेव्हा
$\Sat{M'}{\varphi}$.
\end{proof}

\begin{prob}
\ref{mod:bas:iso:thm:isom} मधील (b) चा पुरावा सविस्तर पूर्ण करा.
\ref{mod:bas:iso:defn:isomorphism} मध्ये समरूपणाचे वैशिष्ट्य ठरवणाऱ्या
पाचही गुणधर्मांपैकी प्रत्येक गुणधर्म कुठे वापरला आहे, ते नोंदवण्याची
खात्री करा.
\end{prob}

\begin{defn}
रचना $\Struct{M}$ ची \emph{स्वयंरूपता} म्हणजे $\Struct{M}$ पासून
त्याच रचनेवर असलेले समरूपण होय.
\end{defn}

\begin{prob}
कोणतीही रचना $\Struct{M}$ असू द्या. जर $X$ हा $\Struct{M}$ चा परिभाष्य
उपसंच असेल आणि $h$ ही $\Struct{M}$ ची स्वयंरूपता असेल, तर $X
= \Setabs{h(x)}{x \in X}$ हे दाखवा (म्हणजे $h$ अंतर्गत $X$ अपरिवर्तित
राहतो).
\end{prob}













\section{संरचनेची उपपत्ती}\label{mod:bas:thm:sec}

प्रत्येक संरचना~$\Struct{M}$ काही वाक्ये सत्य ठरवते आणि
काही असत्य ठरवते. ती सत्य ठरवत असलेल्या सर्व वाक्यांच्या संचाला
तिची \emph{उपपत्ती} म्हणतात. तो संच खरोखरच एक उपपत्ती असतो, कारण
त्याच्यापासून तार्किकरीत्या निष्पन्न होणारे काहीही त्याच्या सर्व
प्रतिमानांमध्ये, आणि त्यांत~$\Struct{M}$ चाही समावेश होतो, सत्य असले
पाहिजे.

\begin{defn}
  संरचना~$\Struct M$ दिलेली असताना, $\Struct{M}$ ची
  \emph{उपपत्ती} म्हणजे $\Struct{M}$ मध्ये सत्य असलेल्या वाक्ये
  चा संच $\Theory{M}$ होय; म्हणजे, $\Theory{M} =
  \Setabs{\varphi}{\Sat{M}{\varphi}}$.
\end{defn}

अभिप्रेत अर्थनिर्धारण असलेल्या वाक्यांच्या संचांसाठीही आपण
``उपपत्ती'' हा शब्द अनौपचारिकपणे वापरतो, मग ते संच निगामीदृष्ट्या बंद
असोत किंवा नसोत.

\begin{prop}
कोणत्याही $\Struct{M}$ साठी $\Theory{M}$ संपूर्ण असते.
\end{prop}

\begin{proof}
कोणत्याही वाक्य~$\varphi$ साठी एकतर $\Sat{M}{\varphi}$ किंवा
$\Sat{M}{\lnot \varphi}$; म्हणून एकतर $\varphi \in \Theory{M}$ किंवा
$\lnot \varphi \in \Theory{M}$.
\end{proof}

\begin{prop}\label{mod:bas:thm:prop:equiv}
  प्रत्येक $\varphi \in \Theory{M}$ साठी $\Struct{N} \models \varphi$ असेल,
  तर $\Struct{M} \elemequiv \Struct{N}$.
\end{prop}

\begin{proof}
सर्व $\varphi \in \Theory{M}$ साठी $\Sat{N}{\varphi}$ असल्यामुळे
$\Theory{M} \subseteq \Theory{N}$. जर $\Sat{N}{\varphi}$, तर
$\Sat/{N}{\lnot \varphi}$; म्हणून $\lnot \varphi \notin \Theory{M}$.
$\Theory{M}$ संपूर्ण असल्यामुळे $\varphi \in \Theory{M}$. म्हणून
$\Theory{N} \subseteq \Theory{M}$ आणि आपल्याला
$\Struct{M} \elemequiv \Struct{N}$ मिळते.
\end{proof}

\begin{rem}\label{mod:bas:thm:remark:R}
  $\Struct{R} = \langle\Real, <\rangle$ विचारात घ्या. या
  संरचनेचा प्रांत म्हणजे वास्तव संख्यांचा संच~$\Real$ असून,
  ती फक्त एक 2-स्थानी विधेय असलेल्या भाषा साठीची रचना
  आहे; त्या विधेयाचे अर्थनिर्धारण वास्तव संख्यांवरील $<$ संबंध असे केले
  आहे. स्पष्टपणे $\Struct{R}$ ही अगणनीय आहे; पण
  $\Theory{R}$ उघडपणे सुसंगत असल्यामुळे लोव्हेनहाइम--स्कोलेम प्रमेयाने
  तिला एक गणनीय प्रतिमान असते, त्याला $\Struct{S}$ म्हणू.
  \ref{mod:bas:thm:prop:equiv} नुसार $\Struct{R} \equiv \Struct{S}$. शिवाय,
  $\Struct{R}$ आणि $\Struct{S}$ या समरूपी नसल्यामुळे यावरून
  \ref{mod:bas:iso:thm:isom} चा उलट दावा सर्वसाधारणपणे खोटा ठरतो.
\end{rem}











\section{आंशिक समरूपणे}\label{mod:bas:pis:sec}

\begin{defn}
  दोन संरचना $\Struct{M}$ आणि $\Struct{N}$ दिलेल्या असताना,
  $\Struct{M}$ पासून $\Struct{N}$ कडेचे \emph{आंशिक समरूपण} म्हणजे
  $\Domain M$ मधील घटक आदान म्हणून घेऊन $\Domain N$ मध्ये मूल्ये देणारे
  असे सांत अंशतः फलन $p$, जे आपल्या प्रांतावर
  \ref{mod:bas:iso:defn:isomorphism} मधील समरूपणाच्या अटी पूर्ण करते:
  \begin{enumerate}
  \item $p$ हे एकास-एक आहे;
  \item प्रत्येक स्थिरांक~$c$ साठी: $p(\Assign{c}{M})$ परिभाषित
    असेल, तर $p(\Assign{c}{M}) = \Assign{c}{N}$;
  \item प्रत्येक $n$-स्थानी विधेय $P$ साठी: $a_1$, \dots, $a_n$
    हे $p$ च्या प्रांतात असतील, तर $\langle a_1, \dots, a_n\rangle \in
    \Assign P M$ तेव्हा आणि केवळ तेव्हाच जेव्हा $\langle p(a_1), \dots, p(a_n) \rangle
    \in \Assign P N$;
  \item प्रत्येक $n$-स्थानी फलन $f$ साठी: $a_1$, \dots, $a_n$
    हे $p$ च्या प्रांतात असतील, तर
    \begin{multline*}
    p(\Assign f M (a_1, \dots,a_n))\\
    = \Assign f N (p(a_1), \dots, p(a_n)).
    \end{multline*}
  \end{enumerate}
  $p$ सांत असण्याचा अर्थ $\dom{p}$ सांत असतो.
\end{defn}

लक्षात घ्या की रिक्त फलन~$\emptyset$ हे कोणत्याही दोन संरचना
मधील आंशिक समरूपण नेहमीच असते.

\begin{defn}\label{mod:bas:pis:defn:partialisom}
  दोन संरचना $\Struct{M}$ आणि $\Struct{N}$ यांना
  \emph{आंशिकरीत्या समरूपी} म्हणतात आणि $\Struct{M} \iso[p]
  \Struct{N}$ असे लिहितात, तेव्हा आणि केवळ तेव्हाच जेव्हा
  $\Struct{M}$ आणि $\Struct{N}$ यांमधील आंशिक समरूपणांचा अरिक्त संच $I$
  पुढील \emph{पुढे-मागे} गुणधर्म पूर्ण करतो:
  \begin{enumerate}
  \item (\emph{पुढे}) प्रत्येक $p \in I$ आणि $a \in \Domain M$ साठी असा
    $q \in I$ असतो की $p \subseteq q$ आणि $a$ हा $q$ च्या प्रांतात
    असतो;
  \item (\emph{मागे}) प्रत्येक $p \in I$ आणि $b \in \Domain N$ साठी असा
    $q \in I$ असतो की $p \subseteq q$ आणि $b$ हा $q$ च्या व्याप्तीत
    असतो.
  \end{enumerate}
\end{defn}

\begin{thm}\label{mod:bas:pis:thm:p-isom1}
  $\Struct{M} \iso[p] \Struct{N}$ असून $\Struct{M}$ आणि
  $\Struct{N}$ या गणनीय असतील, तर $\Struct{M} \iso
  \Struct{N}$.
\end{thm}

\begin{proof}
  $\Struct{M}$ आणि $\Struct{N}$ या गणनीय असल्यामुळे
  $\Domain{M} = \{a_0, a_1, \ldots \}$ आणि $\Domain{N} = \{b_0, b_1,
  \ldots \}$ असे मानू. मनमाना $p_0 \in I$ घेऊन सुरुवात करून, आंशिक
  समरूपणांची वाढती क्रमिका $p_0 \subseteq p_1 \subseteq p_2
  \subseteq \cdots$ पुढीलप्रमाणे परिभाषित करतो:
  \begin{enumerate}
  \item $n+1$ विषम असेल, म्हणजे $n = 2r$ असेल, तर पुढे गुणधर्म वापरून
    असा $p_{n+1} \in I$ शोधा की $p_n \subseteq p_{n+1}$
    आणि $a_r$ हा $p_{n+1}$ च्या प्रांतात असेल;
  \item $n+1$ सम असेल, म्हणजे $n+1 =2r$ असेल, तर मागे गुणधर्म वापरून
    असा $p_{n+1} \in I$ शोधा की $p_n \subseteq p_{n+1}$
    आणि $b_r$ हा $p_{n+1}$ च्या व्याप्तीत असेल.
  \end{enumerate}
आता
\[p = \bigcup_{n\ge 0} p_n,
\]
असे ठेवले, तर $p$ हे $\Struct{M}$ आणि $\Struct{N}$ यांमधील समरूपण
असते.
\end{proof}

\begin{prob}
  \ref{mod:bas:pis:thm:p-isom1} मध्ये परिभाषित केलेले $p$ खरोखरच
  समरूपण आहे, हे सविस्तर दाखवा.
\end{prob}

\begin{thm}\label{mod:bas:pis:thm:p-isom2}
  $\Struct{M}$ आणि $\Struct{N}$ या निव्वळ संबंधात्मक भाषा
  साठीच्या संरचना आहेत असे समजा (म्हणजे फक्त विधेये
  असलेली आणि फलने किंवा स्थिरे नसलेली भाषा). मग
  $\Struct{M} \iso[p] \Struct{N}$ असेल, तर
  $\Struct{M} \elemequiv \Struct{N}$ सुद्धा असते.
\end{thm}

\begin{proof}
  सूत्रे वरील विगमनाने पुढील गोष्ट दाखवता येते. $a_1$, \dots,
  $a_n$ आणि $b_1$, \dots, $b_n$ यांसाठी प्रत्येक $a_i$ ला $b_i$ कडे
  नेणारे आंशिक समरूपण $p$ असेल, तसेच $s_1(x_i) =a_i$ आणि $s_2(x_i) =b_i$
  ($i =1$, \dots,~$n$ साठी) असेल, तर $\Sat{M}{\varphi}[s_1]$ तेव्हा आणि
  केवळ तेव्हाच जेव्हा $\Sat{N}{\varphi}[s_2]$. $n=0$ चे प्रकरण
  $\Struct{M} \elemequiv \Struct{N}$ देते.
\end{proof}

\begin{rem}
फलने उपस्थित असतील, तरी मागील निष्कर्ष सत्य असतो; पण $a_1$,
\dots, $a_n$ यांनी निर्मित $\Struct{M}$ च्या subसंरचना आणि
$b_1$, \dots, $b_n$ यांनी निर्मित $\Struct{N}$ च्या
subसंरचना यांमध्ये $p$ ने प्रेरित केलेले समरूपण विचारात घ्यावे
लागते.
\end{rem}

मागील निष्कर्षाचे टप्प्यांत ``विघटन'' करता येते: सूत्रामध्ये
एकमेकांत अंतर्भूत संख्यापकांची संख्या आणि संबंधित आंशिक समरूपणांचा
किती वेळा विस्तार करता येतो यांचा संबंध प्रस्थापित करावा लागतो.

\begin{defn}
  कोणतेही सूत्र~$\varphi$ दिले असता, $\varphi$ चा
  \emph{संख्यापक-प्रतांक}, $\QuantRank{\varphi} \in \Nat$, म्हणजे $\varphi$ मध्ये एकमेकांत अंतर्भूत
  संख्यापकांची कमाल संख्या; तो पुनरावर्ती रीतीने परिभाषित केला जातो.
  दोन संरचना $\Struct{M}$ आणि $\Struct{N}$ यांना
  \emph{$n$-सममूल्य} म्हणतात आणि $\Struct{M} \elemequiv[n] \Struct{N}$
  असे लिहितात, जर संख्यापक-प्रतांक जास्तीत जास्त~$n$ असलेल्या सर्व
  वाक्यांबाबत त्या एकमत असतील.
\end{defn}

\begin{prop}\label{mod:bas:pis:prop:qr-finite}
  $\Lang{L}$ ही सांत निव्वळ संबंधात्मक भाषा असू द्या, म्हणजे
  सांत इतकी विधेये आणि स्थिरांक असलेली, पण फलने
  नसलेली भाषा. मग प्रत्येक $n \in \Nat$ साठी, तार्किक
  सममूल्यतेपर्यंत, भाषा $\Lang{L}$ मध्ये संख्यापक-प्रतांक
  जास्तीत जास्त $n$ असलेली प्रथम-क्रम वाक्ये सांत इतकीच असतात.
\end{prop}

\begin{proof}
  $n$ वरील विगमनाने.
\end{proof}

\begin{defn}
  संरचना~$\Struct{M}$ दिलेली असताना, $\Domain M^{<\omega}$ हा
  $\Domain{M}$ वरील सर्व सांत क्रमिकांचा संच असू द्या. घटकांच्या सांत
  क्रमिकांवर $\mathbf{a}, \mathbf{b}, \mathbf{c}, \ldots$ ही चले वापरतो.
  जर $\mathbf{a} \in \Domain{M}^{<\omega}$ आणि $a \in \Domain{M}$,
  तर $\mathbf{a}a$ हे $\mathbf{a}$ आणि $a$ यांची \emph{जोडणी} दर्शवते.
\end{defn}

\begin{defn}
  संरचना $\Struct{M}$ आणि $\Struct{N}$ दिलेल्या असताना, समान
  लांबीच्या क्रमिकांमधील संबंध $I_n \subseteq \Domain M^{<\omega}
  \times \Domain N^{<\omega}$ हे $n$ वरील पुनरावर्तनाने पुढीलप्रमाणे
  परिभाषित करतो:
   \begin{enumerate}
   \item $I_0(\mathbf{a},\mathbf{b})$ तेव्हा आणि केवळ तेव्हाच जेव्हा
     $\mathbf{a}$ आणि $\mathbf{b}$ या अनुक्रमे $\Struct{M}$ आणि
     $\Struct{N}$ मध्ये तीच आण्विक सूत्रे पूर्ण करतात; म्हणजे,
     त्या क्रमिकांची समान लांबी k असेल, $s_1(x_i) = a_i$ आणि $s_2(x_i) =
     b_i$ असेल आणि $\varphi$ हे आण्विक असून त्यातील सर्व चले
     $x_1$, \dots,~$x_k$ यांपैकी असतील, तर $\Sat{M}{\varphi}[s_1]$ तेव्हा आणि
     केवळ तेव्हाच जेव्हा~$\Sat{N}{\varphi}[s_2]$.

   \item $I_{n+1} (\mathbf{a},\mathbf{b})$ तेव्हा आणि केवळ तेव्हाच
     जेव्हा प्रत्येक $a\in \Domain M$ साठी असा $b\in \Domain N$ असतो की
     $I_n (\mathbf{a}a,\mathbf{b}b)$, आणि उलटही.
   \end{enumerate}
\end{defn}


\begin{defn}
  $\Struct{M}$ आणि $\Struct{N}$ यांच्यासाठी
  $I_n(\emptyseq,\emptyseq)$ खरे असेल, तर $\Struct{M} \approx_n
  \Struct{N}$ असे लिहितो (येथे $\emptyseq$ ही रिक्त क्रमिका आहे).
\end{defn}

\begin{thm}\label{mod:bas:pis:thm:b-n-f}
  $\Lang{L}$ ही निव्वळ संबंधात्मक भाषा असू द्या. मग
  $I_n (\mathbf{a},\mathbf{b})$ असल्यास, $\QuantRank{\varphi} \le n$ असलेल्या
  प्रत्येक $\varphi$ साठी $\Sat{M}{\varphi}[\mathbf{a}]$ तेव्हा आणि केवळ तेव्हाच
  जेव्हा $\Sat{N}{\varphi}[\mathbf{b}]$ (येथे पुन्हा, $s(x_i) = a_i$ असणारे
  कोणतेही $s$ जर $\varphi$ पूर्ण करत असेल, तर $\mathbf{a}$ हे $\varphi$ पूर्ण
  करते). शिवाय, $\Lang{L}$ सांत असेल, तर उलट दावाही खरा असतो.
\end{thm}

\begin{proof}
  $I_n(\mathbf{a},\mathbf{b})$ यावरून $\mathbf{a}$ आणि $\mathbf{b}$
  संख्यापक-प्रतांक जास्तीत जास्त $n$ असलेली तीच सूत्रे पूर्ण करतात,
  हा पुरावा $\varphi$ वरील सोप्या विगमनाने मिळतो. उलट दाव्यासाठी आपण $n$
  वरील विगमन करतो आणि \ref{mod:bas:pis:prop:qr-finite} वापरतो; त्यावरून प्रत्येक
  $n$ साठी त्या संख्यापक-प्रतांकाची परस्पर तार्किकदृष्ट्या असममूल्य
  सूत्रे जास्तीत जास्त सांत इतकी असतात.

  $n=0$ साठी, $\mathbf{a}$ आणि $\mathbf{b}$ ही तीच संख्यापकमुक्त
  सूत्रे पूर्ण करतात या गृहीतकावरून ती तीच आण्विक सूत्रे पूर्ण करतात;
  म्हणून $I_0(\mathbf{a},\mathbf{b})$.

  $n+1$ च्या प्रकरणासाठी, $\mathbf{a}$ आणि $\mathbf{b}$ ही
  संख्यापक-प्रतांक जास्तीत जास्त $n+1$ असलेली तीच सूत्रे पूर्ण
  करतात असे समजा. $I_{n+1}(\mathbf{a},\mathbf{b})$ दाखवण्यासाठी प्रत्येक
  $a \in \Domain M$ साठी असा $b \in \Domain N$ असतो की
  $I_n(\mathbf{a}a,\mathbf{b}b)$, एवढे दाखवणे पुरेसे आहे; आणि विगमन
  गृहीतकाने प्रत्येक $a \in \Domain M$ साठी असा $b \in \Domain N$ असतो
  की $\mathbf{a}a$ आणि $\mathbf{b}b$ ही संख्यापक-प्रतांक जास्तीत जास्त
  $n$ असलेली तीच सूत्रे पूर्ण करतात, एवढे दाखवणे पुन्हा पुरेसे आहे.

  $a \in \Domain M$ दिलेला असताना, $\Struct{M}$ मध्ये $\mathbf{a}a$ ने
  पूर्ण केलेल्या, प्रतांक जास्तीत जास्त $n$ असलेल्या
  $\psi(x,\mathbf{y})$ या सूत्रांचा संच $\mathsf{T}^a_n$ असू द्या.
  $\mathsf{T}^a_n$ सांत आहे, म्हणून त्यातील सूत्रांचा संयोग असलेले एकच प्रथम-क्रम
  सूत्र म्हणून त्याचा विचार करू शकतो. यावरून $\mathbf{a}$ हे
  $\lexists[x][\mathsf{T}^a_n(x,\mathbf{y})]$ पूर्ण करते; या सूत्राचा
  संख्यापक-प्रतांक जास्तीत जास्त $n+1$ आहे. गृहीतकानुसार $\mathbf{b}$ हे
  $\Struct{N}$ मध्ये तेच सूत्र पूर्ण करते; म्हणून असा $b \in
  \Domain N$ असतो की $\mathbf{b}b$ हे $\mathsf{T}^a_n$ पूर्ण करते. विशेषतः,
  $\mathbf{b}b$ ही $\mathbf{a}a$ प्रमाणेच संख्यापक-प्रतांक जास्तीत जास्त
  $n$ असलेली सर्व सूत्रे पूर्ण करते. त्याचप्रमाणे, प्रत्येक $b \in
  \Domain N$ साठी असा $a\in \Domain M$ असतो की $\mathbf{a}a$ आणि
  $\mathbf{b}b$ ही संख्यापक-प्रतांक जास्तीत जास्त $n$ असलेली तीच
  सूत्रे पूर्ण करतात, हे दाखवता येते; त्यामुळे पुरावा पूर्ण होतो.
\end{proof}

\begin{cor}\label{mod:bas:pis:cor:b-n-f}
  $\Struct{M}$ आणि $\Struct{N}$ या सांत भाषा मधील निव्वळ
  संबंधात्मक संरचना असतील, तर $\Struct{M} \approx_n\Struct{N}$
  तेव्हा आणि केवळ तेव्हाच जेव्हा $\Struct{M} \elemequiv[n] \Struct{N}$.
  विशेषतः, $\Struct{M} \elemequiv \Struct{N}$ तेव्हा आणि केवळ तेव्हाच
  जेव्हा प्रत्येक $n$ साठी $\Struct{M} \approx_n \Struct{N}$.
\end{cor}











\section{सघन रेषीय क्रम}\label{mod:bas:dlo:sec}

\begin{defn}
  \emph{अंतबिंदुविरहित सघन रेषीय क्रम} म्हणजे एकच 2-स्थानी
  विधेय~$<$ असलेल्या भाषा साठीची आणि पुढील वाक्ये पूर्ण
  करणारी संरचना~$\Struct{M}$ होय:
  \begin{enumerate}
  \item $\lforall[x][\lnot x < x]$;
  \item $\lforall[x][\lforall[y][\lforall[z][(x < y \lif (y < z \lif x
    <z ))]]]$;
  \item $\lforall[x][\lforall[y][(x< y \lor \eq[x][y] \lor y < x)]]$;
  \item $\lforall[x][\lexists[y][x < y]]$;
  \item $\lforall[x][\lexists[y][y < x]]$;
  \item $\lforall[x][\lforall[y][(x < y \lif \lexists[z][(x < z \land
        z < y))]]]$.
 \end{enumerate}
\end{defn}

\begin{thm}\label{mod:bas:dlo:thm:cantorQ}
  कोणतेही दोन गणनीय अंतबिंदुविरहित सघन रेषीय क्रम समरूपी
  असतात.
\end{thm}

\begin{proof}
  $\Struct{M_1}$ आणि $\Struct{M_2}$ हे गणनीय अंतबिंदुविरहित
  सघन रेषीय क्रम असू द्या, जेथे ${<_1} = \Assign{<}{M_1}$ आणि ${<_2} =
  \Assign{<}{M_2}$; तसेच $\PIso{I}$ हा त्यांच्यामधील सर्व आंशिक
  समरूपणांचा संच असू द्या. किमान $\emptyset \in \PIso{I}$ असल्यामुळे
  $\PIso{I}$ रिक्त नाही. $\PIso{I}$ हा पुढे-मागे गुणधर्म पूर्ण करतो हे
  आपण दाखवतो. मग $\Struct{M_1} \iso[p] \Struct{M_2}$; आणि
  \ref{mod:bas:pis:thm:p-isom1} नुसार प्रमेय सिद्ध होते.

  $\PIso{I}$ हा पुढे गुणधर्म पूर्ण करतो हे दाखवण्यासाठी $p \in
  \PIso{I}$ असू द्या, $i = 1$, \dots,~$n$ साठी $p(a_i) = b_i$ असू द्या
  आणि सर्वसाधारणतेला बाधा न आणता $a_1 <_1 a_2 <_1 \cdots <_1 a_n$
  असे समजा. $a \in \Domain{M_1}$ दिलेला असताना, तो आधीच त्या नकाशाच्या
  प्रांतात असेल, तर तोच नकाशा अपेक्षित विस्तार आहे. नकाशा रिक्त असेल,
  तर दुसऱ्या रचनेच्या अरिक्त प्रांतातील कोणताही घटक प्रतिमा म्हणून
  निवडून त्याची दिलेल्या घटकाबरोबरची जोडी त्या नकाशात घालता येते.
  उरलेल्या प्रकरणांत
  पुढीलप्रमाणे $b \in \Domain{M_2}$ शोधा:

  \begin{enumerate}
  \item जर $a <_1 a_1$, तर $b <_2 b_1$ असेल असा $b \in
    \Domain{M_2}$ घ्या;
  \item जर $a_n <_1 a$, तर $b_n <_2 b$ असेल असा $b \in \Domain{M_2}$ घ्या;
 \item एखाद्या $i$ साठी $a_i <_1 a <_1 a_{i+1}$ असेल, तर $b_i <_2 b
   <_2 b_{i+1}$ असेल असा $b \in \Domain{M_2}$ घ्या.
  \end{enumerate}
  $\Struct{M_2}$ हा अंतबिंदुविरहित सघन रेषीय क्रम असल्यामुळे अपेक्षित
  गुणधर्म असलेला $b$ नेहमी शोधता येतो. $q = p \cup \{ \langle a, b
  \rangle \}$ असे परिभाषित करा; मग $q \in \PIso{I}$ हा $p$ चा अपेक्षित
  विस्तार आहे. यावरून पुढे गुणधर्म प्रस्थापित होतो. मागे गुणधर्मही
  याचप्रमाणे सिद्ध होतो. म्हणून $\Struct{M_1} \iso[p]
  \Struct{M_2}$; \ref{mod:bas:pis:thm:p-isom1} नुसार $\Struct{M_1} \iso
  \Struct{M_2}$.
\end{proof}

\begin{prob}
  $\PIso{I}$ हा मागे गुणधर्म पूर्ण करतो, हे तपासून
  \ref{mod:bas:dlo:thm:cantorQ} चा पुरावा पूर्ण करा.
\end{prob}

\begin{rem}
  $\Struct{S}$ हा कोणताही गणनीय अंतबिंदुविरहित सघन रेषीय क्रम
  असू द्या. मग (\ref{mod:bas:dlo:thm:cantorQ} नुसार) $\Struct{S} \iso
  \Struct{Q}$, जेथे $\Struct{Q} = (\Rat, <)$ हा परिमेय संख्यांचा संच
  $\Rat$ प्रांत म्हणून असलेला गणनीय सघन रेषीय क्रम आहे. आता
  \ref{mod:bas:thm:remark:R} मधील संरचना~$\Struct{R} = (\Real, <)$
  पुन्हा विचारात घ्या. असा गणनीय संरचना~$\Struct{S}$
  आहे की $\Struct{R} \elemequiv \Struct{S}$, हे आपण पाहिले. पण
  $\Struct{S}$ हा गणनीय अंतबिंदुविरहित सघन रेषीय क्रम आहे;
  म्हणून तो संरचना~$\Struct{Q}$ शी समरूपी (आणि म्हणून प्राथमिक
  सममूल्य) आहे. प्राथमिक सममूल्यतेच्या संक्रमणकतेने $\Struct{R}
  \elemequiv \Struct{Q}$. (हाच पुढे-मागे युक्तिवाद वापरून
  $\Struct{R} \iso[p] \Struct{Q}$ प्रस्थापित करून आपण हे थेटही दाखवू
  शकलो असतो.)
\end{rem}


